\subsection{局部可数分割}

定义 7. --- 称子集所成的一个集族 \(\mathfrak{A}\)（拓扑空间 \(T\) 上的）是局部可数的，如果对每个 \(t \in T\)，存在一个邻域 \(V\)（\(t\) 的），使得那些 \(A \in \mathfrak{A}\)（与 \(V\) 相交者）所成之集是可数的。

若 T 的子集所成之集 A 是局部可数的，则对 T 的每个紧子集 K，与 K 相交的那些 \(A \in A\) 所成之集是可数的，因为 K 可被点的开邻域中的有限个所覆盖，而其中每个开邻域只与 A 中属于 A 的可数个子集相交。

定义 7 表明，局部紧空间的局部可数的 \(\mu\)-可测（相应地，局部 \(\mu\)-可忽略）子集所成集族的并是 \(\mu\)-可测（相应地，局部 \(\mu\)-可忽略）的（No.~1, 命题 3 与 No.~2, 命题 5）。

命题 14. --- 设 X 是一个局部紧空间，\(\mu\) 是 X 上的一个测度，A 是 X 的一个 \(\mu\)-可测子集，\(\mathfrak{K}\) 是在 A 中 \(\mu\)-稠密的 A 的一个紧子集族。存在一个局部可数的集族 \(\mathfrak{H} \subset \mathfrak{K}\)，由两两不相交的集合组成，使得 \(A - \bigcup_{K \in \mathfrak{H}} K\) 局部 \(\mu\)-可忽略，且对每个 \(K \in \mathfrak{H}\)，\(\mu_K\) 的支集恰为整个 K。

考虑由两两不相交的集合组成的诸集族 \(\mathfrak{L} \subset \mathfrak{K}\)，使得对每个 \(L \in \mathfrak{L}\)，\(\operatorname{Supp}(\mu_L) = L\) 。这些集族 \(\mathfrak{L}\) 构成一个非空子集 \(\mathscr{H}\)（\(\mathfrak{P}(\mathfrak{K})\) 中的，因为它含有元素 \(\varnothing\) ），且将按 \(\mathfrak{P}(\mathfrak{K})\) 中的包含关系排序。直接看出 \(\mathscr{H}\) 是归纳的；设 \(\mathfrak{H}\) 是 \(\mathscr{H}\) 的一个极大元素（S, R, §6, No.~10）。首先证明 \(\mathfrak{H}\) 是局部可数的。事实上，对每个 \(x \in X\)，设 \(V\) 是 \(x\) 的一个相对紧的开邻域；若 \((K_i)_{1 \leqslant i \leqslant n}\) 是 \(\mathfrak{H}\) 中与 \(V\) 相交的集合所成的有限族，则

\[
\sum_ {i = 1} ^ {n} | \mu | \left(\mathrm{K} _ {i} \cap \mathrm{V}\right) = | \mu | \left(\mathrm{V} \cap \left(\bigcup_ {i = 1} ^ {n} \mathrm{K} _ {i}\right)\right)
\]

因为诸 \(K_{i}\) 两两不相交，故 \(\sum_{i=1}^{n}| \mu |(\mathrm{K}_{i} \cap \mathrm{V}) \leqslant |\mu |(\mathrm{V})\) 。由此推出，若 \(H_{V}\) 就是那些 \(K \in \mathfrak{H}\)（与 V 相交者）所成之集，则

\[
\sum_ {\mathrm{K} \in \mathfrak {H} _ {\mathrm{V}}} | \mu | (\mathrm{K} \cap \mathrm{V}) <   + \infty ,
\]

且由于 \(|\mu|(\mathrm{K} \cap \mathrm{V}) > 0\) 对每个 \(\mathrm{K} \in \mathfrak{H}_{\mathrm{V}}\) 成立，故 \(\mathfrak{H}_{\mathrm{V}}\) 必为可数。其次，证明 \(\mathrm{N} = \mathrm{A} - \bigcup_{\mathrm{K} \in \mathfrak{H}} \mathrm{K}\) 是局部 \(\mu\)-可忽略的。我们在上面已看到 \(\mathrm{N}\) 是 \(\mu\)-可测的。若 \(\mathrm{N}\) 不是局部可忽略的，它就包含一个非可忽略的紧集 \(\mathrm{L}_0\)，从而（No.~8, 命题 12）包含一个非可忽略紧集 \(\mathrm{L} \subset \mathrm{L}_0\)（属于 \(\mathfrak{K}\)）。由于 \(|\mu_{\mathrm{L}}|(\mathrm{L}) = |\mu|(\mathrm{L}) > 0\)（No.~7, 引理 2 与引理 3），测度 \(\mu_{\mathrm{L}}\)（在 \(\mathrm{L}\) 上由 \(\mu\) 诱导）非零；其支集 \(\mathrm{S}\) 因而是属于 \(\mathfrak{K}\) 的非空紧集（由 \((\mathrm{PL}_{\mathrm{I}})\)），且 \(\operatorname{Supp}(\mu_{\mathrm{S}}) = \mathrm{S}\)（No.~7, 引理 2, (iii)）。由此推出，集族 \(\mathfrak{H} \cup \{\mathrm{S}\}\) 属于 \(\mathscr{H}\)，这与 \(\mathfrak{H}\) 的定义矛盾；于是集 \(\mathrm{N}\) 是局部可忽略的，证毕。

